% Generated draft: requires mathematical and source review.
\documentclass[11pt]{article}
\usepackage[margin=25mm]{geometry}
\usepackage{fontspec}
\setmainfont{Latin Modern Roman}
\usepackage{amsmath,amssymb,mathtools,unicode-math}
\setmathfont{Latin Modern Math}
\usepackage{enumitem,needspace,xurl,hyperref}
\setlength{\parindent}{0pt}
\setlength{\parskip}{5pt}
\setlength{\emergencystretch}{3em}
\newcommand{\N}{\mathbb{N}}
\newcommand{\Z}{\mathbb{Z}}
\newcommand{\Q}{\mathbb{Q}}
\newcommand{\R}{\mathbb{R}}
\newcommand{\C}{\mathbb{C}}
\newcommand{\F}{\mathbb{F}}
\newcommand{\A}{\mathbb{A}}
\newcommand{\PP}{\mathbb{P}}
\newcommand{\E}{\mathbb{E}}
\newcommand{\eps}{\varepsilon}
\newcommand{\dd}{\,\mathrm{d}}
\begin{document}
\section*{UnsolvedMath ID: 30000511}
\textbf{Source problem code:} catalogue entry
\subsection{Definitions and mathematical statement}
Determine which moduli spaces of algebraic curves are rational.

For an integral finite-type variety $V$ over a field $k$, rational over $k$ means birational over $k$ to a projective space $\PP^r_k$, where $r=\dim V$. Equivalently,
\[
 k(V)\cong k(t_1,\ldots,t_r)
 \quad\text{as extensions of }k.
\]
Rationality is stronger than merely being dominated by a rational variety (unirationality), and it is not replaced here by stable rationality or rational connectedness.

The source asks for a classification of the intended curve-moduli spaces according to this property, not a conjecture that every moduli space is rational. It gives no ground field, genus, number of marked points, level or other auxiliary structure, curve class, or parameter range. Nor does it identify coarse moduli varieties versus moduli stacks, possible twists, or the irreducible components to which variety rationality should apply. Ordinary birational rationality of a coarse integral variety and an appropriate rationality notion for a stack are different targets.

Consequently the entry remains a broad historical rationality-classification request. Before a precise problem can be stated, the report must identify the moduli family and the field over which rationality is tested. No selection of a familiar family such as unpointed smooth-curve moduli, a pointed family, or genus-zero forms is silently imposed, and no list of rational or nonrational genera is invented.
\subsection{Short English statement}
The source asks for a rationality classification of curve-moduli spaces, but does not identify the spaces or base field. Rationality retains its birational meaning and is not weakened to unirationality or another nearby property.
\subsection{Sources}
\begin{itemize}
\item Local UnsolvedMath registry, permanent ID 30000511, code catalogue entry, ``Rationality of Moduli Spaces of Algebraic Curves.'' Proposed year 2006; proposer unspecified.
\item \emph{Classical Algebraic Geometry} (2007), Oberwolfach Reports, pp.\ 1615--1662, DOI 10.4171/owr/2006/27; \url{https://doi.org/10.4171/owr/2006/27}. The displayed report year 2007 differs from the DOI and proposed year 2006.
\item Supplied reference ``Florence--Hoffmann--Reichstein, rationality problems'': \url{https://arxiv.org/abs/1709.05696}. The precise moduli family covered is not given in this record.
\item Statement status exact; the supplied verification says official-report content was checked and notation, context, and question form normalized without changing the claim.
\item Literature triage checked 2026-08-21: open. It describes the sentence as a historical fragment without fixed parameters and says the general rationality classification remains unresolved. This does not supply a status for any chosen individual moduli family. Created 2026-08-25; updated 2026-09-16. Time-sensitive claims require independent verification.
\item Registry curation by Ulam AI / UnsolvedMath Contributors, CC BY 4.0; underlying report and literature provenance remain as supplied. Unreviewed statement-only draft.
\end{itemize}
\end{document}
